\subsection{Local stability and convergence}
\label{sec:local-stability-convergence}

\subsubsection{Local stability and retained curvature}
\label{sec:local-stability}

With fixed mobility, the pullback changes the local stability condition at a
stationary point \(\theta_*\) of \(F\). Throughout this subsection a
subscript \(*\) denotes evaluation at that point, so that
\(H_*=\nabla^2F(\theta_*)\) and \(B_*=B_\alpha(\theta_*)\), and
\((\theta_*,p_*,q_*)\) is the corresponding state of
\cref{eq:abstract-momentum-step,eq:abstract-momentum-states}. The comparison
depends on the Hessian curvature left after subtracting twice \(B_*\).

\begin{theorem}
\label{thm:local-schur-stability}
Let \(\theta_*\) satisfy \(\nabla F(\theta_*)=0\), suppose the component
energies are twice continuously differentiable near \(\theta_*\) with
positive shifted energies, and let \(H_*\succ0\). Fix
\(0\leq\alpha\leq1\), \(M\succ0\), \(\eta>0\), and
\(0\leq\beta<1\). The corresponding rest state has \(p_*=0\) and a stored
scalar \(q_*>0\); set \(\kappa_*=q_*/Q(\theta_*)\). The parameter--momentum block of the
linearization of \cref{eq:abstract-momentum-step,eq:abstract-momentum-states}
is Schur stable, meaning that all its eigenvalues have modulus less than one,
if and only if
\begin{equation}
  \kappa_*H_*
  \prec 2B_*+\frac{2(1+\beta)}{\eta}M^{-1}.
  \label{eq:local-stability-criterion}
\end{equation}
\end{theorem}

The proof, in \cref{app:schur-proof}, reduces the linearized
parameter--momentum map to a matrix polynomial and then to a scalar quadratic
in the Rayleigh quotients \(b=\xi^*\widehat B_*\xi/(\xi^*\xi)\) and
\(h=\xi^*\widehat H_*\xi/(\xi^*\xi)\), where
\(\widehat B_*=M^{1/2}B_*M^{1/2}\), \(\widehat H_*=M^{1/2}H_*M^{1/2}\),
and \(\xi^*\) is the conjugate transpose of \(\xi\in\mathbb C^d\).

The unrelaxed scalar equation has a neutral direction at the stationary state;
the scalar is controlled separately in the local convergence argument below.

\begin{corollary}
\label{cor:local-step-interval}
Under the assumptions of \cref{thm:local-schur-stability}, define
\begin{equation}
  L_{\mathrm{eff}}
  =\lambda_{\max}\!\left(
  M^{1/2}(\kappa_*H_*-2B_*)M^{1/2}\right).
  \label{eq:effective-local-stiffness}
\end{equation}
If \(L_{\mathrm{eff}}>0\), the exact stability interval is
\begin{equation}
  0<\eta<\frac{2(1+\beta)}{L_{\mathrm{eff}}}.
  \label{eq:exact-local-step-interval}
\end{equation}
If \(L_{\mathrm{eff}}\leq0\), every finite \(\eta>0\) is admissible.
At the stationary point,
\begin{equation}
  B_1(\theta_*)=0,
  \qquad B_*=\alpha B_m(\theta_*),
  \qquad \rank B_m(\theta_*)\leq\min\{d,m-1\}.
  \label{eq:stationary-component-curvature}
\end{equation}
For a common \(\theta_*\), \(q_*\), \(M\), and \(\beta\), increasing
\(\alpha\) cannot shrink the stability interval relative to the aggregate
choice \(\alpha=0\).
\end{corollary}

\begin{proof}
Apply a congruence by \(M^{1/2}\) to
\cref{eq:local-stability-criterion}. The total gradient vanishes at
\(\theta_*\), giving \(B_1(\theta_*)=0\). The component gradients sum
to zero and are therefore linearly dependent, which gives the rank bound.
Finally, \(B_m\succeq0\) makes \(L_{\mathrm{eff}}\) nonincreasing with
\(\alpha\).
\end{proof}

\begin{remark}
\label{rem:stationary-cancellation}
Gradient cancellation can leave nonzero component curvature after the
aggregate correction has vanished; if all component gradients vanish, so does
the component correction. Since \(H_*\succ0\), stability for every
\(\eta>0\) requires \(B_*\succ0\), which \cref{eq:stationary-component-curvature}
excludes once \(m-1<d\); a low-rank correction can nevertheless enlarge a
finite interval. The comparison in \cref{cor:local-step-interval} is between
component and aggregate pullbacks within the momentum update;
\cref{app:appended-comparison} records the corresponding comparison with the
frozen appended recurrence, whose stability interval is at least as large.
\end{remark}

\subsubsection{Local convergence under any scalar relaxation}
\label{sec:local-relaxed-convergence}

\begin{theorem}
\label{thm:local-relaxed-convergence}
Let the component energies be twice continuously differentiable near
\(\theta_*\), with positive shifted energies,
\(\nabla F(\theta_*)=0\), and \(H_*\succ0\). Fix
\(M\succ0\), \(\eta>0\), \(0\leq\beta<1\), and
\(0\leq\alpha\leq1\), and assume
\begin{equation}
  H_*\prec2B_\alpha(\theta_*)+
  \frac{2(1+\beta)}{\eta}M^{-1}.
  \label{eq:local-relaxed-stability-condition}
\end{equation}
Use the momentum update with \(M^n\equiv M\), the relaxation
\cref{eq:generic-relaxation} with any sequence \(\rho^n\in[0,1]\), and
\(q^0=Q(\theta^0)\). There exist \(\delta>0\), \(c_0,c_1<\infty\),
and \(r\in(0,1)\) such that, whenever
\(a_0=\norm{\theta^0-\theta_*}+\norm{p^0}<\delta\), the iterates remain
near \(\theta_*\) and
\begin{equation}
  \norm{\theta^n-\theta_*}+\norm{p^n}\leq c_0r^na_0.
  \label{eq:local-geometric-convergence}
\end{equation}
Moreover, \(q^n\to q_\infty>0\), with
\begin{equation}
  0\leq Q(\theta_*)-q_\infty\leq c_1a_0^2.
  \label{eq:local-scalar-limit}
\end{equation}
The constants are uniform over all such relaxation sequences, including
\(\rho^n\equiv1\).
\end{theorem}

The proof in \cref{app:local-convergence-proof} combines contraction of the
parameter--momentum map with a quadratic bound on the scalar drift.
